\section{Reproducibility and Traceability}

This section specifies the minimal reproducibility and traceability guarantees
required of an OPS-1 Engine Core artefact. These guarantees are \emph{structural}:
they concern the ability to reproduce diagnostic results from declared inputs
only. They do not concern the empirical behavior of the diagnosed system.

\subsection{Reproducibility as architectural invariant (binary)}

OPS-1 treats reproducibility as a \textbf{binary architectural property}:
an artefact is either reproducible or it is not. No intermediate status,
gradation, or confidence level is admissible.

\paragraph{Necessary condition for validity.}
Given identical declared inputs, the Engine Core must reproduce identical
diagnostic results. If this condition fails, the artefact is invalid under OPS-1.

\paragraph{Not a process property.}
Reproducibility is not a matter of process discipline, execution hygiene,
toolchain choice, or environment management. It is fixed exclusively at the
level of kernel semantics and declared inputs.

\paragraph{Non-reproducibility implies terminal STOP-admissibility.}
If reproducibility cannot be established for an output, diagnostic admissibility
is violated. The situation is terminal STOP-admissible, and the output must not
be treated as an OPS-1 result.

\subsection{Registry and evidence as architecture (non-technical)}

OPS-1 treats the \textbf{append-only registry} and \textbf{immutable evidence} as
architectural constructs, not as tooling, storage systems, or processes.

\paragraph{Registry (logical, reference-only).}
The registry is a logical reference space that contains \emph{references} and
\emph{links} only. It carries no interpretation, no aggregation semantics, and no
decision or filter logic. Registry contents must not influence evaluation,
classification, terminal STOP semantics, or admissibility boundaries.

\paragraph{Evidence (immutable, context-free).}
Evidence denotes immutable artefacts referenced by the registry. Evidence is
context-free in OPS-1 kernel semantics: within OPS-1 the Engine Core must not
interpret, normalize, clean, sample, filter, aggregate, or otherwise transform
evidence.

Evidence is immutable after creation: it must never be overwritten, replaced, or
deleted. Any change to content constitutes a different artefact and must be
represented only as a new, separately referencable evidence artefact.

\subsection{Append-only invariants}

\paragraph{Append-only registry rule.}
Registry entries are append-only. Once a reference or link exists, it must not be
mutated, replaced, reordered, or deleted. Any registry that permits mutation is
non-admissible under OPS-1.

\paragraph{Order-stable registry semantics.}
Registry introduction order is semantically stable in the logical sense that
references, once introduced, remain introduced. No re-ordering is permitted to
change meaning, provenance identity, or diagnostic interpretation.

\paragraph{Append-only evidence referencing.}
Once an evidence artefact is referenced, that reference must remain valid without
post-hoc alteration. Any later artefact with changed content must be treated as a
distinct evidence artefact and referenced separately.

\subsection{Provenance (complete or invalid)}

Each diagnostic result $R$ must be traceable to its declared inputs via a
\textbf{complete} provenance mapping:
\[
(S,\mathcal{I},\Theta) \;\mapsto\; R
\]

\paragraph{Completeness requirement.}
Provenance is \emph{complete or invalid}. There is no admissible notion of partial,
probabilistic, approximate, inferred, or reconstructed provenance within OPS-1.
If any required link is missing, provenance is invalid.

\paragraph{Transitive traceability.}
Every diagnostic output must be transitively traceable to:
\begin{itemize}
  \item the specific input snapshot $S$;
  \item the specific invariant set $\mathcal{I}$;
  \item the specific threshold objects $\Theta$;
  \item the Engine Core version identifier;
  \item all evidence artefacts referenced for that output (if any).
\end{itemize}

The traceability requirement concerns \emph{referencing} and \emph{link integrity}
only. It provides no explanatory narrative, causal account, or justification.

\subsection{Stable identifiers and manifests (format-agnostic)}

The Engine Core must expose stable identifiers for:
\begin{itemize}
  \item input state snapshots;
  \item invariant sets;
  \item threshold objects;
  \item the Engine Core version.
\end{itemize}

The choice of identifier algorithms, encodings, manifests, or serialization
formats is external to kernel semantics, provided identifiers are stable,
inspectable, and unambiguous.

\subsection{Explicit prohibitions (non-negotiable)}

The following are explicitly non-admissible under OPS-1:

\begin{itemize}
  \item \textbf{Mutable registry:} any ability to edit, rewrite, replace, reorder,
        or delete registry entries.
  \item \textbf{Evidence normalization or cleaning:} any transformation of
        evidence content that changes meaning, regardless of labeling.
  \item \textbf{Aggregation with semantic effect:} any aggregation that alters
        the meaning of evidence or changes diagnostic interpretation.
  \item \textbf{Post-hoc correction of evidence:} any after-the-fact repair,
        replacement, rewriting, or retrofitting of evidence to enforce consistency.
  \item \textbf{Opaque preprocessing:} any uninspectable transformation or hidden
        intermediate that breaks traceability.
  \item \textbf{Hidden sampling or filtering:} any undisclosed selection,
        exclusion, prioritization, or weighting that prevents complete provenance.
  \item \textbf{Semantically relevant non-determinism:} any non-determinism that
        can influence diagnosis, terminal STOP evaluation, or provenance identity.
\end{itemize}

These prohibitions are architectural invariants. They are independent of
implementation strategy and must hold for OPS-1 conformance.
